\section{Mean Bias-Aware Low-Bit Training Method}
\label{sec:lowbit_pilot}
\label{sec:lowbit_method}

Low-bit training amplifies the pathological effects of anisotropy.
Under blockwise quantization, scale is determined by extreme elementwise values.
As shown in Section~\ref{sec:activation_outliers},
these extremes are predominantly mean-dominated.
This suggests that addressing mean bias at its source may substitute
for heavier spectral corrections.


We test the following hypothesis:
\begin{quote}
\emph{If mean bias is the dominant driver of outliers,
then source-level mean removal should recover most of the stability
benefits attributed to SVD-based spike suppression methods such as Metis.}
\end{quote}

We propose a simple \emph{mean--residual quantization} method, which we call
\emph{Averis} (Averaging-Induced Residual Splitting).
The method targets a concrete empirical observation: for large activation-like matrices in Transformer
training, extreme values (which dominate blockwise quantization scales) are largely driven by a
mean-dominated component that is coherent across tokens.
Accordingly, we explicitly split out a column-mean component via averaging and quantize it separately
from the remaining residual, using only reduction operations and standard quantization kernels.

\paragraph{Notation and quantization operator.}
Let $\mathcal{Q}_b(\cdot)$ denote a $b$-bit quantization operator (e.g., blockwise FP4) applied to a matrix.
We write $\bar{\mathbf{M}}\triangleq \mathcal{Q}_b(\mathbf{M})$ for the quantized version of any matrix
$\mathbf{M}$.
Throughout, $\mathbf{1}$ denotes an all-ones vector with the appropriate dimension.

\paragraph{Forward pass: activation mean--residual splitting.}
Let $\mathbf{W}\in\mathbb{R}^{m\times n}$ be a weight matrix and let
$\mathcal{X}\in\mathbb{R}^{b\times s\times m}$ be an input activation tensor, where $b$ is batch size,
$s$ is sequence length, and $m$ is hidden dimension.
As usual, we reshape $\mathcal{X}$ into a matrix $\mathbf{X}\in\mathbb{R}^{l\times m}$ with $l=b\cdot s$
and compute the GeMM output
\[
\mathbf{Y}=\mathbf{X}\mathbf{W}\in\mathbb{R}^{l\times n}.
\]

We define the \emph{activation column-mean vector} $\mu_{\mathbf{X}}\in\mathbb{R}^{m}$ by
\[
\mu_{\mathbf{X}} \triangleq \frac{1}{l}\mathbf{1}^{\top}\mathbf{X},
\]
The residual is
\[
\mathbf{X}_{\mathrm{R}} \triangleq \mathbf{X}-\mathbf{1}\,\mu_{\mathbf{X}}.
\]
We then quantize the mean vector and residual independently:
\[
\bar{\mu}_{\mathbf{X}}=\mathcal{Q}_b(\mu_{\mathbf{X}}),
\qquad
\bar{\mathbf{X}}_{\mathrm{R}}=\mathcal{Q}_b(\mathbf{X}_{\mathrm{R}}),
\qquad
\bar{\mathbf{W}}=\mathcal{Q}_b(\mathbf{W}).
\]
The quantized forward computation is
\begin{equation}
\label{eq:averis_forward}
\hat{\mathbf{Y}}
\;=\;
\mathbf{1}\,(\bar{\mu}_{\mathbf{X}}\bar{\mathbf{W}})
\;+
\bar{\mathbf{X}}_{\mathrm{R}}\,\bar{\mathbf{W}}.
\end{equation}
where the first term is broadcast along the token dimension via $\mathbf{1}$.
In practice, one computes $\mu_{\mathbf{X}}$ first, subtracts it from $\mathbf{X}$ on the fly to form
$\mathbf{X}_{\mathrm{R}}$, and never materializes the mean matrix.

\paragraph{Backward pass: output-gradient mean--residual splitting.}
Let $\mathbf{D}\in\mathbb{R}^{l\times n}$ denote the matrix of loss derivatives with respect to
the GeMM output $\mathbf{Y}$, i.e.,
\[
\mathbf{D}\triangleq \frac{\partial \mathcal{L}}{\partial \mathbf{Y}}
\quad (\text{after reshaping}).
\]
We define the \emph{output-gradient column-mean vector} $\mu_{\mathbf{D}}\in\mathbb{R}^{n}$ by
\[
\mu_{\mathbf{D}} \triangleq \frac{1}{l}\mathbf{1}^{\top}\mathbf{D},
\qquad
\mathbf{D}_{\mathrm{R}} \triangleq \mathbf{D}-\mathbf{1}\,\mu_{\mathbf{D}}.
\]
Quantize the components independently:
\[
\bar{\mu}_{\mathbf{D}}=\mathcal{Q}_b(\mu_{\mathbf{D}}),
\qquad
\bar{\mathbf{D}}_{\mathrm{R}}=\mathcal{Q}_b(\mathbf{D}_{\mathrm{R}}).
\]
Then the key GeMMs in backpropagation are evaluated as
\begin{equation}
\label{eq:averis_backward}
\widehat{\frac{\partial \mathcal{L}}{\partial \mathbf{X}}}
\;=\;
\mathbf{1}\,(\bar{\mu}_{\mathbf{D}}\bar{\mathbf{W}}^{\top})
\;+
\bar{\mathbf{D}}_{\mathrm{R}}\,\bar{\mathbf{W}}^{\top},
\qquad
\widehat{\frac{\partial \mathcal{L}}{\partial \mathbf{W}}}
\;=\;
\bar{\mathbf{X}}_{\mathrm{R}}^{\top}\bar{\mathbf{D}}_{\mathrm{R}}
\;+
\bar{\mathbf{X}}_{\mathrm{R}}^{\top}(\mathbf{1}\bar{\mu}_{\mathbf{D}})
\;+
(\mathbf{1}\bar{\mu}_{\mathbf{X}})^{\top}\bar{\mathbf{D}}_{\mathrm{R}}
\;+
(\mathbf{1}\bar{\mu}_{\mathbf{X}})^{\top}(\mathbf{1}\bar{\mu}_{\mathbf{D}}).
\end{equation}
The additive terms are accumulated with broadcasted additions without explicitly storing
$\mathbf{M}_{\mathbf{X}}$ or $\mathbf{M}_{\mathbf{D}}$.

\paragraph{Computational profile.}
Averis requires only computing two mean vectors ($\mu_{\mathbf{X}}$ and $\mu_{\mathbf{D}}$) and applying
two subtractions per step, in addition to the baseline quantized GeMMs.
No spectral decomposition or iterative orthogonalization is involved, and all additional operations are
GPU-friendly reductions and elementwise kernels.

\section{Mean Bias-Aware Low-Bit Training Experiments}
\label{sec:lowbit_experiment}

\subsection{Experimental Setup}

\label{experiment:setup}
\textbf{Models and Datasets.}
We conduct experiments on Qwen-3 (0.6B), and we additionally include a Qwen3-8B pretraining loss reference run.
For pretraining, we use the DCLM dataset and train the 0.6B model on 100B tokens.
The raw data are segmented to split long documents and concatenate short ones, filtered to remove non-Unicode and non-English content.

\textbf{FP4 Training.}
We evaluate the efficacy of Averis under FP4 quantization against BF16 baselines.
All FP4 training in this work adopts W4A4G4 quantization, where weights, activations, and gradients are represented in the E2M1 NVFP4 format.
Stochastic rounding (SR) is applied by default in all FP4 experiments, as it mitigates quantization bias and is orthogonal to Averis.

\textbf{Baselines.}
We compare the following settings:
\begin{itemize}[leftmargin=*,itemsep=2pt]
\item \textbf{BF16}: full-precision baseline.
\item \textbf{Vanilla FP4}: direct NVFP4 quantization (W4A4G4) without structural splitting.
\item \textbf{Averis FP4}: the proposed mean--residual splitting applied to activations and output gradients, with independent quantization of mean vectors and residuals.
\end{itemize}

\subsection{Results}

\textbf{Training loss.}
Figure~\ref{figure:loss-averis} shows the training loss curves of the BF16 baseline and FP4 + Averis.
Averis remains slightly above BF16 in loss, while clearly improving over vanilla FP4.

\begin{figure}[h]
    \centering
    \includegraphics[width=0.92\linewidth]{figures/qwen_0.6B_loss.png}
    \caption{Training loss curves for Qwen3-0.6B under BF16 and FP4 + Averis.}
    \label{figure:loss-averis}
\end{figure}

\textbf{Qwen3-8B pretraining loss.}
Figure~\ref{figure:loss-averis-8b} presents the corresponding pretraining loss result on Qwen3-8B under NVFP4 training.

\begin{figure}[h]
    \centering
    \includegraphics[width=0.92\linewidth]{figures/8B-nvfp4-loss.png}
    \caption{Pretraining loss result for Qwen3-8B under NVFP4 training.}
    \label{figure:loss-averis-8b}
\end{figure}

\textbf{Downstream performance at 10B tokens.}
At the 10B-token checkpoint, we evaluate BF16 and Averis on seven downstream tasks.
As shown in Table~\ref{table:downstream-10b-averis}, Averis improves the average score
from 0.4564 to 0.4661.

\begin{table}[h]
\centering
\caption{Downstream task results at 10B tokens.}
\label{table:downstream-10b-averis}
\small
\begin{tabular}{lcccccccc}
\toprule
Method & ARC-C & ARC-E & BoolQ & HellaSwag & LAMBADA & PIQA & RACE & Avg \\
\midrule
BF16        & \textbf{0.2534} & \textbf{0.5126} & 0.5309 & \textbf{0.3768} & 0.3602 & \textbf{0.6730} & 0.4882 & 0.4564 \\
Averis      & 0.2491 & 0.5072 & \textbf{0.5746} & 0.3751 & \textbf{0.3862} & 0.6670 & \textbf{0.5036} & \textbf{0.4661} \\
\bottomrule
\end{tabular}
\end{table}

% \subsection{Computational Cost}

% We report the computational overhead of Averis relative to Vanilla FP4 and BF16 baselines.
% \textbf{Placeholder}: Throughput (tokens/s), step time (ms), and relative overhead (\%) at matched global batch size.

% \subsection{Memory Usage}

% We report training memory consumption for all compared methods.
% \textbf{Placeholder}: Peak GPU memory (GB), activation memory (GB), and optimizer-state memory (GB)
% under identical model size and training configuration.
